\section{Related Work}
\label{sec:related}

\subsection{Vision-Based Social Navigation}
Traditional and learning-based social navigation (SocialNav) systems optimize paths to prevent physical collisions and respect implicit social norms~\cite{PrinciplesGuidelinesEvaluating2025}. Trajectory forecasting approaches proactively yield space based on future pedestrian states~\cite{CognitionPrecognitionFutureaware2025}, while recent foundation models distill human-inspired social compliance into unified navigation policies~\cite{SocialNavTrainingHumanInspired2026} or employ vision-language models (VLMs) to evaluate path etiquette~\cite{VLMsocialnavSociallyAware2025, ObstaclesEtiquetteRobot2026}. However, these vision-centric policies inherently depend on LOS sensors. Unable to perceive NLOS risks, they fail to anticipate a dynamic pedestrian approaching from a blind corner or yield to unobserved human activities, often causing abrupt boundary violations when entities suddenly fall into LOS.

\subsection{Audio-Visual Navigation}
AVNav equips agents with omnidirectional auditory perception to overcome restricted visual fields. Initial research proved sound effectively guides agents to static acoustic targets in unexplored spaces~\cite{LookListenAct2020, LearningSetWaypoints2021}. Beyond continuous signals, recent frameworks localize transient semantic events~\cite{SemanticAudiovisualNavigation2021}, track moving sources~\cite{CatchMeIf2023}, and tackle multi-source interference~\cite{AudiovisualNavigationNoisy2025}. Advanced architectures also introduce interactive conversational agents to query human oracles~\cite{AVLENAudiovisuallanguageEmbodied2022, CAVENEmbodiedConversational2024} or integrate large language models (LLMs) for semantic reasoning and instruction-following~\cite{RILAReflectiveImaginative2024}. However, these approaches treat non-target sounds exclusively as alternative goals or background noise to be attenuated. Consequently, they fail to model social sound events as spatial boundaries, leaving agents prone to disturbing stationary conversations or colliding with NLOS dynamic pedestrians while over-optimizing to reach the acoustic target.

\subsection{Multimodal Spatial Mapping and Memory}
Connecting acoustic signals with spatial reasoning relies on robust multi-modal perception. Foundation models like CLAP~\cite{CLAPLearningAudio2023} and ImageBind~\cite{ImageBindOneEmbedding2023} provide rich zero-shot cross-modal alignments, driving the evolution of open-vocabulary SELD systems~\cite{OpenvocabularySoundEvent2025}. Meanwhile, purely acoustic SELD has scaled through large-scale pre-training~\cite{PSELDNetsPretrainedNeural2025} and joint distance estimation~\cite{ExperimentalStudyJoint2025}, while audio-visual variants leverage visual cues to suppress reverberations~\cite{MVANetMultistageVideo2025}. Furthermore, spatial audio capabilities have been directly injected into foundation models to enable acoustic scene reasoning~\cite{BATLearningReason2024} and sound object segmentation~\cite{SAMAudioSegment2025}. Despite their perceptual capabilities, these feature extractors primarily focus on instantaneous event analysis rather than projecting sounds into persistent spatial constraints for navigation.

To support long-horizon tasks, spatial maps aggregate sensory features into structured representations, such as open-vocabulary fields and 3D scene graphs~\cite{3DDynamicScene2020, ConceptFusionOpensetMultimodal2023, ConceptGraphsOpenvocabulary3D2024, ClioRealtimeTaskdriven2024, BareQueriesOpenvocabulary2025, M3DMapObjectawareMultimodal2025, MultimodalSpatialLanguage2025}. However, purely visual maps cannot register entities before they enter LOS. To mitigate this and handle dynamic changes, embodied memory models update spatial states online~\cite{DynamemOnlineDynamic2025} or maintain acoustic memory maps for setting navigation waypoints~\cite{LearningSetWaypoints2021}. More closely related to our goal, AVLMaps~\cite{MultimodalSpatialLanguage2025} fuse audio into semantic grids. However, they naively register sound events at the agent's current location when a decibel threshold is exceeded, effectively degenerating omnidirectional NLOS acoustics into strictly proximity-based perceptions under static environment assumptions. Similarly, probability-based scene graphs update regional beliefs from coarse audio directionality but stringently require predefined human-activity heatmaps to resolve spatial ambiguities~\cite{HomeEmergencyUsingAudio2025}.

By contrast, our mapping module, SAVMap, seamlessly integrates these perceptual and mapping paradigms while eliminating their restrictions. It anchors transient, unconstrained sound events into the topological layout as inferred entities, and projects these non-visual inferences into a continuous social cost field. This effectively transforms multidirectional acoustic cues into persistent, geometry-aware boundaries for proactive planning.
